\subsection*{BR-012: two unique permanent extensions fail the coupled incidence}
\textbf{Status: proved.  Independent peer whole-proof and source review
passed, including the core intersection, quotient, and root modules.
A fresh peer standard-library reconstruction passed in $0.2014$ seconds.
The earlier complete flag computations establishing uniqueness remain
explicit dependencies of BR-010, whose family ranks were not rerun.}

Let
\[
 S_1=\{0,1,2,3,4,5,8,9,12,13\},\qquad
 S_2=\{0,1,2,3,4,5,7,10,12,13\}.
\]
Their masks are $13119$ and $13503$.  Both literal permanent residual
tensors have border rank at least eighteen in characteristic zero.

\paragraph{The previously complete unique-extension reduction.}
Both tensors have dimensions $10\times15\times15$, with used variables
$x_0,\ldots,x_{14}$ in row-major order.  In the ordered slice basis
$a_0,\ldots,a_9$, the only invariant seventeen-plane remaining after
the complete BR-010 flag and pencil tests is
\[
 E_1=T_1(C^*)+\langle a_7\otimes x_0,a_7\otimes x_3\rangle,
 \qquad
 E_2=T_2(C^*)+\langle a_7\otimes x_0,a_7\otimes x_2\rangle.
\]
The two complete source records contain respectively $65$ and $71$
projective families, all excluded including their finite continuations.
Each has two terminal rigid candidates before the stronger rigid test;
one is excluded by the fourth $B$-Koszul map, of rank respectively
$17043$ and $17057$, with threshold $17\binom{14}{4}=17017$.
The one displayed seventeen-plane remains, with no unhandled family.
The new checker validates this complete structural index cover and hashes
the source records; it does not claim to rerun those previous family ranks.

The common connected solvable group has the three opposite row
translations $r_i\mapsto r_i+s r_0$ in $B$ and
$r_i\mapsto r_i-s r_0$ in $C$, for $i=1,2,3$, together with the common
row operations $r_2\mapsto r_2+s r_1$ and
$r_3\mapsto r_3+s r_1$, and the full row/column torus.  Its action
on $A$ is the previously verified negative dual action on the slice span.
Every opposite root has zero action on $A$, verified coefficient by
coefficient.  Thus switching $B$ and $C$ leaves a common root unchanged
and multiplies an opposite root by $-1$.  The invariant $AB$ and $AC$
subspaces are consequently the same.  In particular the unique admissible
$AC$ plane is the same displayed $E_i$ with $B$ replaced by $C$.

\paragraph{The coupled necessity, including all intersection dimensions.}
For a hypothetical border-rank-seventeen tensor, use the closure of the
incidence of seventeen general rank-one triples together with their
seventeen-dimensional spans in $ABC$, $AB$, $AC$, and $BC$.  In the
generic incidence all three pairwise spans and the triple span have
dimension seventeen.  Its closure in the product of the corresponding
Grassmannians is proper over the tensor line, and its image contains the
full seventeenth secant variety.  Thus the tensor has a nonempty complete
fiber.  This fiber is invariant under the connected solvable stabilizer,
so the fixed-point theorem gives a simultaneous invariant tuple
\[
 (D,E_{AB},E_{AC},E_{BC}),\qquad\dim D=\dim E_{AB}=\dim E_{AC}=\dim E_{BC}=17.
\]
The relevant containment identities hold on the generic incidence and
are closed, so they hold on its closure.  In particular
\[
 T\in D\subset(E_{AB}\otimes C)\cap(E_{AC}\otimes B)
                 \cap(A\otimes E_{BC}).
\]
The pairwise spans lie in the same rank-one Grassmannian limit loci used
for the BR-010 extension tests.  Their uniqueness therefore forces
$E_{AB}=E_{AC}=E_i$.

Put
\[
 V=(E_i\otimes C)\cap(E_i\otimes B),\qquad
 R=\langle A^*\mathbin{\lrcorner}V\rangle\subset B\otimes C,
 \qquad U=T_i(A^*).
\]
Set $G=E_{BC}\cap R$.  Then $G$ is torus invariant,
\[
 U\subset G\subset R,\quad 10\le\dim G\le17,\quad
 \dim\bigl(V\cap(A\otimes G)\bigr)\ge17.
\]
The last inclusion follows because $V\subset A\otimes R$, so $D$ lies
in $A\otimes(E_{BC}\cap R)$.  This argument allows \emph{every}
dimension from ten through seventeen.  It does not enlarge $G$ to a
seventeen-plane or assume that such enlargement preserves a rank bound.

\paragraph{Projection-dimension inequalities.}
For each tensor the exact rational calculation gives
\[
 \dim V=27,\qquad\dim R=36,\qquad\dim U=10.
\]
The quotient $R/U$ has twenty torus weight spaces $Q_0,\ldots,Q_{19}$,
of dimensions
\[
 (\underbrace{1,\ldots,1}_{10},
  \underbrace{2,\ldots,2}_{6},
  \underbrace{1,\ldots,1}_{4}).
\]
Let $d_j=\dim(G/U)_j$.  Consequently $d_j\ge0$, each $d_j$ is bounded
by the displayed weight-space dimension, and
$\sum_{j=0}^{19}d_j\le7$.  These constraints include all possible
torus-fixed $G$, even without imposing the additional root invariance or
border-rank conditions on $G$.

The space $V$ has seventeen weights; write
$W_s=V_s\cap(A\otimes G)$ and $z_s=\dim W_s$.  A projection from
$V_s$ to one quotient weight $Q_j$ involves a unique $A$ coordinate,
because the ten $A$ weights are distinct.  If that projection has rank
$r_{sj}$, its inverse image of $(G/U)_j$ has dimension at most
\[
 \dim V_s-r_{sj}+d_j.
\]
Also $z_s\le\dim V_s$.  Intersecting all required inverse images can
only decrease dimension.  The following complete upper bounds needed
for the proof result from the independently recomputed projection ranks.
The exact torus weights, quotient bases, and all projection ranks are
stored in the certificate; indices in this table use that deterministic
quotient order and the lexicographic order on triple weights.
\[
\begin{array}{c|l|l}
s& S_1\text{ upper bound for }z_s&S_2\text{ upper bound for }z_s\\\hline
0&1+\min_{0\le j<10}d_j&1+\min_{0\le j<10}d_j\\
1&\min_{10\le j<16}d_j&\min_{10\le j<16}d_j\\
2&d_{16}&d_{16}\\
3&d_{16}&d_{17}\\
4&d_{12}&d_{12}\\
5&d_{17}&d_{16}\\
6&d_{11}&2+\min(d_{14},d_{11})\\
7&2+\min(d_{10},d_{15})&d_{15}\\
8&d_{18}&d_{19}\\
9&d_{10}&d_{10}\\
10&d_{14}&d_{14}\\
11&d_{13}&d_{18}\\
12&d_{19}&d_{13}\\
13&d_{19}&d_{18}\\
14&\min(2,d_{13},d_{14})&\min(2,d_{13},d_{14})\\
15&d_{19}&d_{18}\\
16&d_{19}&d_{18}
\end{array}
\]

For $S_1$, use the following convex combinations of valid bounds:
\[
\begin{split}
z_0&\le1+d_3,\qquad
z_1\le\tfrac12d_{15}+\tfrac14d_{11}+\tfrac14d_{12},\\
z_7&\le2+\tfrac14d_{10}+\tfrac34d_{15},\qquad
z_{14}\le1+\tfrac14d_{13}+\tfrac14d_{14}.
\end{split}
\]
Together with the other rows of the table, they yield
\[
 \sum_s z_s\le4+d_3+\tfrac54\sum_{j=10}^{15}d_j
                    +2d_{16}+d_{17}+d_{18}+4d_{19}
 \le4+\tfrac54\cdot7+\tfrac34d_{16}+\tfrac{11}{4}d_{19}
 \le\frac{65}{4}<17.
\]
For $S_2$, take instead
\[
 z_0\le1+d_4,\quad
 z_1\le\tfrac14(d_{10}+d_{11}+d_{12}+d_{15}),\quad
 z_6\le2+d_{11},\quad
 z_{14}\le1+\tfrac14d_{13}+\tfrac14d_{14}.
\]
The resulting bound is
\[
 \sum_s z_s\le4+d_4+\tfrac54\sum_{j=10}^{15}d_j
                    +2d_{16}+d_{17}+4d_{18}+d_{19}
 \le4+\tfrac54\cdot7+\tfrac34d_{16}+\tfrac{11}{4}d_{18}
 \le\frac{65}{4}<17.
\]
Both inequalities contradict the coupled necessity.  In particular,
the integer dimension of the intersection is at most sixteen for every
admissible $G$.  This proves both exclusions.

\paragraph{Independent checker and actual computation scope.}
Run \texttt{python scripts/check\_per10\_unique\_coupled.py}.
This script and \texttt{permanent\_coupled\_module.py} use only the
Python standard library and the previously supplied local standard-library
linear algebra.  They reconstruct both tensors, root actions, the two
displayed $E_i$, the intersections $V$, the contraction spans $R$, and
every projection rank.  They verify the two rational weighted inequalities
using exact fractions.  They import no numerical optimizer and no
floating-point arithmetic is used in the checker.

An optional \texttt{--enumerate} mode also enumerates all $253144$
integer dimension vectors satisfying the twenty bounds and total budget
seven.  It checks their corresponding upper bounds individually; both
maxima are sixteen.  The much shorter rational inequality above is
already sufficient, independently of this extra finite enumeration.
An exploratory numerical optimizer was used only to discover candidate
rational weights; those weights were converted to exact fractions and
verified against freshly reconstructed constraints.  No numerical
optimization output is an assumption of the proof.

The complete result is \texttt{certificates/per10\_unique\_coupled.json}.
It explicitly records that the new coupled ranks were recomputed, while
the original BR-010 family ranks were not recomputed by this checker.
The two masks are added by the separate conservative union checker
\texttt{merge\_per10\_current\_gate.py}; the BR-011 search records
remain unchanged.  The independent peer result is preserved in
\texttt{certificates/per10\_coupled\_peer\_rerun.json}.
